%% conclusion:
% - resumen de contribuciones
% - conclusiones principales
% - trabajo futuro

\chapter{Conclusiones y trabajo futuro}
\label{chap:conclusion}

% TODO: sección de conclusiones (pendiente de resultados experimentales)

\section{Trabajo futuro}
\label{sec:trabajo_futuro}

A partir del trabajo desarrollado en este Trabajo de Fin de Grado, se identifican varias líneas de extensión
que podrían abordarse en investigaciones posteriores.

\begin{itemize}

    \item \textbf{Extensión a problemas de optimización multiobjetivo.} En los últimos
          años se observa un crecimiento notable en la literatura relacionada con la aplicación
          de técnicas subrogadas a problemas con múltiples objetivos en conflicto. El
          enfoque desarrollado en este trabajo se ha centrado en optimización de objetivo único,
          lo que deja abierta una posible extensión hacia el ámbito multiobjetivo.

    \item \textbf{Criterios de activación adaptativos para la hibridación \textit{online}.}
          En la hibridación \textit{online} propuesta, la decisión de evaluar una solución
          candidata con el subrogado o con la función real se rige mediante una
          probabilidad fija. Una posible mejora consistiría en sustituir este mecanismo por
          un criterio adaptativo que tome en consideración el estado actual de la búsqueda o
          la fiabilidad estimada del modelo, recurriendo a la evaluación real únicamente cuando
          la predicción del subrogado no ofrezca suficientes garantías. Este tipo de
          enfoque permitiría reducir aún más el número de evaluaciones reales sin comprometer
          la calidad de las soluciones obtenidas.

    \item \textbf{Análisis de la estabilidad temporal de los modelos subrogados.}
          La evaluación \textit{offline} desarrollada permite estudiar la capacidad de los
          modelos para anticipar la etapa inmediatamente posterior de la búsqueda. Como
          ampliación, resultaría de interés evaluar por separado su rendimiento sobre
          horizontes temporales progresivamente más alejados del conjunto de entrenamiento.
          Este análisis permitiría estimar durante cuánto tiempo siguen siendo útiles las
          relaciones aprendidas por el subrogado y determinar con mayor precisión la
          frecuencia de reentrenamiento más adecuada para su integración \textit{online},
          especialmente ante los cambios introducidos por los mecanismos de reinicio.

    \item \textbf{Muestreo inicial mediante \textit{Latin Hypercube Sampling}.} La
          estrategia \textit{online} desarrollada construye el conjunto de entrenamiento del
          subrogado a partir de las soluciones generadas por la propia metaheurística
          durante su ejecución, lo que implica un periodo inicial en el que el modelo no está
          disponible. Una alternativa que permitiría integrarlo desde la primera evaluación
          consistiría en generar un conjunto inicial de puntos mediante \textit{Latin Hypercube
              Sampling} (LHS), una técnica de muestreo que distribuye los puntos de forma más
          representativa que el muestreo aleatorio puro. Estos puntos se evaluarían con la
          función real antes de iniciar la metaheurística, proporcionando al subrogado
          una visión más amplia del espacio de búsqueda desde el comienzo.

\end{itemize}
